% ============================================================
% PREFÁCIO
% ============================================================
\chapter*{Prefácio}
\addcontentsline{toc}{chapter}{Prefácio}

A excelência corporativa não é destino — é **disciplina diária**. Em um mundo marcado por volatilidade extrema (pandemia, guerra, inflação, revolução tecnológica, emergência climática), as organizações que prosperam não são as maiores ou as mais antigas, mas as que **aprendem mais rápido, executam com rigor e adaptam-se com propósito**.

Este livro nasceu da convergência de três fontes: (1) a literatura acadêmica seminal e contemporânea — de Chandler e Porter a Teece, Kaplan, Christensen e Agrawal; (2) a prática gerencial de vanguarda no Brasil — Magazine Luiza, Weg, Nubank, Embraer, Natura, Ambev, JBS, Petrobras, bancos, hospitais; (3) dados reais de 2024–2025 — IBGE, BCB, CVM, B3, CNI, CGI.br, WEF, McKinsey, BCG, Gartner, PwC, Deloitte, KPMG, AON, Protiviti, FNQ, ISO, SBTi, TCFD, ISSB.

A estrutura reflete a **arquitetura da gestão integrada**: 
\begin{itemize}
    \item \textbf{Parte I (Cap. 1–3):} Fundamentos — estratégia, finanças, operações. Onde a organização decide \textit{para onde vai} e \textit{como produz valor}.
    \item \textbf{Parte II (Cap. 4–6):} Pessoas e execução — gente, marketing/vendas, projetos. Onde a estratégia encontra a realidade.
    \item \textbf{Parte III (Cap. 7–9):} Inovação e tecnologia — novos produtos, qualidade, TI/dados. O motor de renovação contínua.
    \item \textbf{Parte IV (Cap. 10–12):} Relacionamento e governança — comunicação/stakeholders, ESG, compliance/governança. A licença para operar e a confiança.
    \item \textbf{Parte V (Cap. 13–14):} Novos horizontes — riscos corporativos, transformação digital \& IA. O preparo para o desconhecido.
\end{itemize}

Cada capítulo segue padrão rigoroso: \textbf{Resumo} (250 palavras), \textbf{Introdução} (contexto, problema, objetivo), \textbf{Fundamentação Teórica} (4–6 eixos com referências seminal + 2020–2025), \textbf{Metodologia} (como aplicamos ao Brasil), \textbf{Análise \& Argumentação} (4 figuras/tabelas originais + 3 cases + dados reais), \textbf{Conclusão} (imperativos gerenciais + agenda pesquisa), \textbf{Referências ABNT}.

Os \textbf{4 artefatos visuais por capítulo} (56 no total) — Strategy Maps, OKR Cascades, Dynamic Capabilities Loops, Blue Ocean Canvases, DuPont Decompositions, Rolling Forecast Timelines, VSMs, Kraljic Matrices, Resilience Dashboards, Logistics Cost Benchmarks, 9-Box Talent Reviews, Leadership Pipelines, Employee Journeys, Engagement Benchmarks, Customer Journey Maps, RevOps Funnels, MarTech Stacks, CAC/LTV Analyses, Hybrid Decision Matrices, Portfolio Kanbans, PMO Dashboards, PPM Maturity Curves, Innovation Funnels, JTBD Maps, Three Horizons Portfolios, Adoption/Chasm Curves, QFD Houses, DMAIC Roadmaps, MEG Scorecards, Cost of Quality Charts, Enterprise Architecture Maps, Data Maturity Models, Cloud Adoption Frameworks, Cyber Risk Heat Maps, Stakeholder Salience Matrices, Crisis Timelines, Employee Journeys, Social Listening Dashboards, Materiality Matrices, Net Zero Roadmaps, Value Chain ESG Maps, Integrated Value Creation Maps, Compliance Program Matrices, Due Diligence Flows, Three Lines Models, Implementation Timelines, Bow-Tie Analyses, Risk Appetite Statements, Corporate Risk Heat Maps, Climate Scenario Maps, AI Maturity Models, GenAI Use Case Portfolios, MLOps Pipelines, Responsible AI Frameworks — são **originais, desenhados em TikZ/pgfplots**, prontos para publicação em vetor.

Este livro não pretende ser "manual de receitas". Pretende ser **referência para decisões difíceis**: onde alocar capital, como estruturar governança, qual risco assumir, quando pivotar, como construir cultura que executa estratégia. A excelência corporativa é a soma de milhares de microdecisões bem tomadas, dia após dia.

Agradeço aos conselheiros, executivos, pesquisadores e estudantes que, direta ou indiretamente, contribuíram com dados, validações e provocações. Erros e omissões são inteiramente meus.

\vspace{1cm}
\noindent\textit{São Carlos, 2026}

\noindent\textbf{Jeronimo Alves dos Santos}
\end{flushright}